
\section{Methodology}

\subsection{Overview}

gPCD exploits the inherent parallelism of the graphics pipeline to perform broad-phase occupancy generation concurrently with point-particle rendering. Each particle is processed independently within the graphics pipeline, allowing the occupancy structure required for collision detection to be constructed while the particle representation is simultaneously rendered.

The method is based on a uniform spatial subdivision in which particles are mapped to a fixed Cartesian cell array. Particle occupancy is determined through corner-based placement of an axis-aligned bounding region (AARB), producing a potentially colliding set (PCS) that is subsequently refined through exact distance-based collision checks.

\subsection{Spatial Cell Array}

The simulation domain is partitioned into a uniform grid of cubic cells of dimension

\[
1 \times 1 \times 1
\]

Each cell is identified by integer coordinates $(i,j,k)$ corresponding to its center location. For memory-efficient addressing, the three-dimensional cell coordinate is transformed into a single linear index,

\[
I = i + N_x \left(j + N_y k\right)
\]

where $N_x$ and $N_y$ are the grid dimensions in the $x$- and $y$-directions. The cell array is therefore stored as a contiguous one-dimensional structure.

Each cell contains a fixed-length particle occupancy list storing the indices of particles wholly or partially contained in that cell. The maximum list length is determined from the largest allowable local occupancy.

\subsection{Particle Representation}

A particle structure array of length equal to the total number of particles is maintained in GPU memory. Each particle stores geometric and kinematic properties together with a corner-location array of length eight.

Each particle is enclosed by an axis-aligned bounding region (AARB) possessing eight corners. These corners are used exclusively for spatial occupancy assignment. A particle diameter constraint is imposed such that

\[
d_p \leq 1
\]

which guarantees that a particle may intersect no more than eight neighboring cells. The worst-case configuration occurs when a particle lies at a common vertex shared by eight cells.

\subsection{Broad-Phase Collision Candidate Generation}

Broad-phase processing is executed concurrently with normal graphics-pipeline activity. Each particle is submitted as a vertex to the vertex shader. For every particle, the vertex stage generates eight corner-specific fragment invocations corresponding to the eight AARB corners.

Each fragment invocation receives a unique corner identifier

\[
c \in \{0,1,\dots,7\}
\]

which maps directly to the associated entry in the particle corner-location array. Because each invocation writes only to its assigned corner slot, write conflicts are avoided without atomic operations.

For each corner, the corresponding cell index is computed and the particle identifier is inserted into that cell's occupancy list. 

The resulting populated cell array defines the potentially colliding set (PCS).

\subsection{Narrow-Phase Collision Detection}

The PCS is transferred directly to the compute pipeline. A compute shader launches one thread per particle. For a source particle $p$, the thread iterates over the eight entries of its corner-location array and retrieves the corresponding occupied cells.

For each candidate particle $q$ listed in those cells, an exact geometric collision test is performed. For spherical particles, contact is determined by

\[
\|\mathbf{x}_p - \mathbf{x}_q\| < r_p + r_q
\]

where $\mathbf{x}$ denotes particle position and $r$ particle radius. Other geometric tests may be substituted for non-spherical particles.

Only candidates contained in shared cells are tested, substantially reducing the number of pairwise checks relative to all-to-all search.

\subsection{Duplicate Elimination}

Two sources of duplicate collision candidates arise naturally:

\begin{enumerate}
    \item Corner duplicates: multiple corners of the same particle occupy the same cell.
    \item Particle duplicates: the same particle pair appears in multiple shared cells.
\end{enumerate}

To remove redundant tests, each compute invocation maintains a local duplicate-cache array storing previously processed target particle identifiers. Before testing a candidate $q$, the cache is queried. If $q$ has already been processed, the candidate is skipped; otherwise, the identifier is inserted and testing proceeds.

This local strategy removes redundant collision checks without requiring inter-thread synchronization or global memory barriers.

\subsection{Computational Characteristics}

Under the particle-size constraint, each particle references at most eight cells, yielding bounded neighbor-search complexity per particle. Let $m$ denote the average number of occupants per visited cell. The expected narrow-phase cost per particle is

\[
\mathcal{O}(8m) = \mathcal{O}(m)
\]

leading to an overall complexity approximately linear in particle count for bounded occupancy distributions.

Because broad-phase occupancy assignment is overlapped with graphics execution, its effective standalone overhead is reduced relative to separate compute-only preprocessing passes.
